\section{Experiments}
\label{sec:experiments}

The main visual experiments are deferred while large-scale pretraining is in
progress. We first test whether the same bidirectional solve operator transfers
outside vision, using genomic sequence classification as an auxiliary
evaluation. This setting is useful because it combines short regulatory motifs
with global sequence context and includes sequence lengths for which quadratic
attention becomes arithmetically expensive.

\subsection{Auxiliary Experiment I: Genomic Sequence Classification}

\paragraph{Datasets and protocol.}
We use eight binary classification tasks from
GenomicBenchmarks~\cite{gresova2023genomic}: coding versus intergenic sequence,
human versus worm, mouse enhancers, the Cohn and Ensembl human enhancer tasks,
human regulatory regions, non-TATA promoters, and human open chromatin regions
(OCR). We retain the official test split and select checkpoints using a
stratified $10\%$ validation split drawn only from the training set. Results are
test-set mean and sample standard deviation over three seeds.

Both mixers use the same compact, non-pretrained DNA encoder: width $256$, two
mixer blocks, eight heads, a learned absolute position embedding, a local
depthwise motif stem of width seven, a $4\times$ MLP, and masked mean pooling.
The \RRLSSO\ model uses rank $r=32$ and fixed one-dimensional rank rotation;
the baseline replaces only the mixer with multi-head attention (MHA). We train
for 100 epochs with AdamW, cosine decay, effective batch size $128$, and BF16.
The shared augmentation applies reverse complementation with probability
$0.5$ and a $0.2\%$ nucleotide mutation rate, followed by 20 clean epochs.
At test time, original and reverse-complement logits are averaged. Thus the
comparison isolates the global mixer while preserving the same local inductive
bias, optimization, augmentation, and pooling.

\twocolumn[{
  \begin{minipage}{\textwidth}
  \centering
  \footnotesize
  \setlength{\tabcolsep}{3.5pt}
  \refstepcounter{table}\label{tab:genomic}
  \textbf{Table \thetable:} GenomicBenchmarks test results and mixer arithmetic.
  Accuracy is reported in percent as mean $\pm$ sample standard deviation over
  three seeds. MACs are per sample for both mixer layers at the configured
  maximum sequence length; one multiply--accumulate is one MAC. Bold marks the
  higher mean.\par\smallskip
  \begin{tabular}{lrrrrrrr}
    \toprule
    Task & $N_{\max}$ & \RRLSSO\ Acc. & MHA Acc. & \RRLSSO\ MCC & MHA MCC
      & \RRLSSO\ GMAC & MHA GMAC \\
    \midrule
    Coding--Intergenic
      & 200  & \textbf{90.03 $\pm$ 0.15} & 89.47 $\pm$ 0.07
      & \textbf{0.8006} & 0.7901 & 0.089 & 0.146 \\
    Human--Worm
      & 200  & \textbf{96.89 $\pm$ 0.10} & 96.73 $\pm$ 0.01
      & \textbf{0.9378} & 0.9345 & 0.089 & 0.146 \\
    Mouse Enhancers
      & 4707 & \textbf{76.03 $\pm$ 0.72} & 73.55 $\pm$ 4.37
      & \textbf{0.5271} & 0.4874 & 2.083 & 25.155 \\
    Cohn Enhancers
      & 500  & \textbf{72.70 $\pm$ 0.55} & 72.28 $\pm$ 0.16
      & \textbf{0.4545} & 0.4474 & 0.222 & 0.518 \\
    Ensembl Enhancers
      & 573  & \textbf{91.50 $\pm$ 0.66} & 90.24 $\pm$ 0.43
      & \textbf{0.8306} & 0.8051 & 0.254 & 0.637 \\
    Regulatory
      & 802  & \textbf{94.11 $\pm$ 0.11} & 94.04 $\pm$ 0.02
      & \textbf{0.9115} & 0.9106 & 0.355 & 1.079 \\
    Non-TATA
      & 251  & 90.71 $\pm$ 1.12 & \textbf{91.07 $\pm$ 0.18}
      & 0.8187 & \textbf{0.8235} & 0.112 & 0.196 \\
    OCR
      & 593  & \textbf{80.33 $\pm$ 0.45} & 79.45 $\pm$ 0.35
      & \textbf{0.6068} & 0.5899 & 0.263 & 0.671 \\
    \midrule
    Macro average
      & -- & \textbf{86.54} & 85.85
      & \textbf{0.7359} & 0.7236 & -- & -- \\
    \bottomrule
  \end{tabular}
  \vspace{1em}
  \end{minipage}
}]

\paragraph{Predictive results.}
As shown in \cref{tab:genomic}, \RRLSSO\ obtains the higher mean accuracy and
MCC on seven of eight tasks. It improves macro-average accuracy by $0.68$
percentage points and MCC by $0.0124$. The only loss is on Non-TATA promoters,
where MHA leads by $0.36$ percentage points. The \RRLSSO\ models also use
exactly 131,552 fewer parameters than their matched MHA counterparts; because
the learned position table changes with task length, this corresponds to a
$4.5\%$--$7.5\%$ reduction in total parameters. These results do not claim a
genomics state of the art; rather, they isolate the transferability of the
proposed mixer under a controlled architecture and training recipe.

\paragraph{Arithmetic and long-sequence behavior.}
For one MHA layer, we count
\begin{equation}
  \operatorname{MAC}_{\mathrm{MHA}}=4ND^2+2N^2D,
\end{equation}
covering the query/key/value and output projections, $QK^\top$, and attention
readout. For \RRLSSO, we count the fused $U/C$ and output projections, compact
statistics, Woodbury readout, and dense-equivalent Cholesky and triangular
solves:
\begin{align}
  \operatorname{MAC}_{\RRLSSO}
  ={}&ND(2D+Hr)+HNr^2+2NrD \notag\\
    &+Dr^2+\frac{Hr^3}{6}.
\end{align}
Elementwise normalization, rotary transforms, activations, and softmax are
excluded from both counts. Across the eight tasks, \RRLSSO\ reduces mixer MACs
by $1.64\times$--$12.08\times$.

Lower arithmetic does not necessarily imply lower wall-clock time for short
sequences, because Flash-SDPA benefits from highly mature, hardware-specific
kernel fusion. The longest task nevertheless demonstrates that the asymptotic
advantage can become practical. Mouse Enhancers has maximum length 4,707 and
mean effective length 2,348. On an NVIDIA GeForce RTX 5070 Ti, using the same
BF16 training step, batch size 64, augmentation, backward pass, gradient
clipping, and fused AdamW update, \RRLSSO\ reduces mixer arithmetic from
$25.16$ to $2.08$ GMAC and shortens one training epoch from $3.18$ to $1.73$
seconds, a $1.83\times$ end-to-end speedup. We therefore treat MACs as the
primary implementation-independent efficiency measure and this longest task as
evidence that the linear sequence dependence can translate into realized
training acceleration.
